\section{The algorithm}\label{ch:algo}
\sectionbrief
{an input in one of the data models M1--M3 (Section~\ref{f:ch:data}) and, for each step, the hypotheses of the theorem that justifies it.}
{to turn the theorems into one procedure with explicit inputs, outputs and guarantees.}
{operational summary of Sections~\ref{f:ch:ident}--\ref{f:ch:nonlin}. Table~\ref{tab:steps} lists the steps, Figure~\ref{fig:algo} gives the control flow, and two worked examples make every step concrete.}

\subsection{The pipeline}

\begin{center}\small
\begin{tabular}{@{}c L{2.6cm} L{5.0cm} L{5.4cm}@{}}
\toprule
Step & Task & Result relied upon & Output\\\midrule
0 & Input & data models M1--M3 & exact $u$\\
1 & Annihilator & Gram matrix of a dictionary & a candidate $P$\\
2 & Factor & unique factorisation (F1) & factors $P_i$, multiplicities $e_i$\\
3 & EP test & Lemma~\ref{f:lem:ep}, Thm.~\ref{f:thm:ident}, Ch.~\ref{f:ch:data} & a verdict per factor\\
4 & Reduce & Thms.~\ref{thm:spec} and~\ref{thm:intrinsic} & $G_u$, $J_u$, certificate\\
5 & Decompose & Thm.~\ref{f:thm:factor}, Eq.~\eqref{f:eq:lagrange}, Thm.~\ref{f:thm:hm} & components $u_j$, ambiguity $\delta(P)$, check $\gamma_N$\\
6 & Terms & Thms.~\ref{f:thm:eint}, \ref{thm:cost}, \ref{f:thm:lie} & log-symbol, $\Eint$\\
7 & Nonlinear branch & Thms.~\ref{f:thm:dict}, \ref{f:thm:jet} & relative identification, flow-based $\Eint$\\
\bottomrule
\end{tabular}
\captionof{table}{The steps of the algorithm: what each step does, which result justifies it, and what it returns.}\label{tab:steps}
\end{center}

\begin{figure}[!htb]
\centering
\begin{tikzpicture}[
  st/.style={draw=cA,fill=cA!8,rounded corners=2pt,minimum height=11mm,text width=5.4cm,align=center,font=\small\sffamily},
  dec/.style={draw=cD,fill=cD!10,diamond,aspect=2.4,inner sep=1pt,align=center,font=\footnotesize\sffamily},
  outb/.style={draw=cC,fill=cC!10,rounded corners=2pt,minimum height=10mm,text width=5.4cm,align=center,font=\small\sffamily},
  arr/.style={-{Stealth[length=2.2mm]},thick,cG}]
\def\dy{1.75}
\node[st] (s0) at (0,0) {\textbf{0. Input}\\ $u$ in model M1, M2 or M3};
\node[st] (s1) at (0,-\dy) {\textbf{1. Annihilator}\\ find $P$ with $P(\partial)u=0$};
\node[st] (s2) at (0,-2*\dy) {\textbf{2. Factor}\\ $P=\prod_iP_i^{e_i}$};
\node[st] (s3) at (0,-3*\dy) {\textbf{3. EP test}\\ is $w_i=\tfrac{P}{P_i}(\partial)u$ in $\EP$?};
\node[dec] (d) at (0,-4.35*\dy) {all tests\\ say $\notin\EP$?};
\node[outb] (o1) at (0,-5.5*\dy) {\textbf{Identified:} $\Ann(u)=(P)$};
\node[st] (s4) at (7.2,-4.35*\dy) {\textbf{4. Reduce}\\ drop factors that test ``$\in\EP$'': $G_u$, $J_u$};
\node[st] (s5) at (7.2,-5.5*\dy) {\textbf{5. Decompose}\\ $u=\sum_ju_j$, check $\gamma_N>0$};
\node[st] (s6) at (7.2,-6.65*\dy) {\textbf{6. Terms}\\ log-symbol, $\Eint$};
\node[st] (s7) at (7.2,-7.8*\dy) {\textbf{7. Nonlinear branch}\\ dictionary, flows of terms};
\draw[arr] (s0)--(s1); \draw[arr] (s1)--(s2); \draw[arr] (s2)--(s3); \draw[arr] (s3)--(d);
\draw[arr] (d)--node[left,font=\footnotesize]{yes}(o1);
\draw[arr] (d)--node[above,font=\footnotesize]{no}(s4);
\draw[arr] (s4)--(s5); \draw[arr] (s5)--(s6); \draw[arr] (s6)--(s7);
\draw[arr,dashed] (o1)--node[above,font=\footnotesize]{then}(s5);
\end{tikzpicture}
\caption{The algorithm as a flowchart. The left column certifies what the operator is. If every factor test says ``not an exponential-polynomial'', the operator is identified. Otherwise Step~4 reduces it to $G_u$. The right column extracts structure from $u$ once the operator is known.}\label{fig:algo}
\end{figure}

\subsection{The steps in detail}

\begin{aside}{Step 1: obtain an annihilator}
Let $\Phi(u)=(T_1u,\dots,T_Mu)^{\top}$ be the vector of $M$ candidate terms applied to $u$, and
\[
\mathbb G(u)=\int_\Omega\Phi\,\Phi^{*}\dd\mu
\]
its Gram matrix (positive semidefinite). A relation $\lambda\cdot\Phi=0$ holds if and only if $\lambda\in\ker\mathbb G(u)$. The computation is exact (rational or algebraic numbers, Gaussian elimination): no learning, no statistics. If $\ker\mathbb G$ has dimension larger than one, $u$ satisfies several relations, which is itself information about $G_u$ and $J_u$.
\end{aside}

\begin{aside}{Step 2: factor}
Factor $P(\tau,\xi)$ into irreducibles with multiplicities. This is a standard computer-algebra operation over the algebraic numbers.
\end{aside}

\begin{aside}{Step 3: exponential-polynomial test}
For each $i$ compute $w_i=(P/P_i)(\partial)u$ and decide whether $\dimc\ring w_i<\infty$. How depends on the data model:
\begin{itemize}
\item M1: a syntactic check (is the closed form a finite sum of polynomials times exponentials?);
\item M3: a holonomic computation, decidable modulo (AC) (Theorem~\ref{thm:holo});
\item M2: undecidable from the jet alone (Theorem~\ref{f:thm:blackbox}); with a bound $B$, compare annihilators up to degree $B+\deg P-1$ (Theorem~\ref{thm:decide}).
\end{itemize}
\end{aside}

\begin{aside}{Step 4: reduce to $G_u$}
Drop from $P$ each factor whose test returned ``EP'', with the multiplicity bookkeeping of Theorem~\ref{thm:intrinsic}: at most $\sum_ie_i$ tests in total. The output is $G_u$; the dropped part is $J_u$, whose finite zero set lists the exponentials contained in $u$. The certificate reads either ``$\Ann(u)=(P)$'' or ``$u$ identifies only $G_u$ and contains these exponentials''.
\end{aside}

\begin{aside}{Step 5: decompose}
For evolution factors $\tau-a_j(\xi)$ apply the Lagrange formula~\eqref{f:eq:lagrange}. For general coprime factors use the construction of Theorem~\ref{f:thm:factor}; the components are determined up to the $\delta(P)$-dimensional exponential-polynomial space, fixed by a normalisation. Before using the components in $H^m$, check $\gamma_N>0$ (Theorem~\ref{f:thm:hm}).
\end{aside}

\begin{aside}{Step 6: terms and interaction}
Inside an evolution factor with symbol $a=\sum_na_n$, the terms are read from the log-symbol, which is linear in time:
\[
\log\hat u(t,\kappa)=\log\hat u_0(\kappa)+t\sum_na_n(i\kappa).
\]
Then evaluate $\Eint$ by the product formula~\eqref{f:eq:eint} at cost $O(N)$ per mode. For non-commuting terms, classify the generated Lie algebra (Theorem~\ref{f:thm:lie}).
\end{aside}

\begin{aside}{Step 7: nonlinear case}
Replace Step~3 by the dictionary criterion (Theorem~\ref{f:thm:dict}) and Step~6 by the flow-based interaction term and its jet (Theorem~\ref{f:thm:jet}).
\end{aside}

\subsection{Worked example 1: the wave equation}

We follow Steps~1--5 on the wave equation with $c=1$ and the solution
\[
u(t,x)=\underbrace{\frac{1}{1+(x+t)^2}}_{f(x+t)}+\underbrace{e^{-(x-t)^2}}_{g(x-t)} .
\]

\begin{example}[{Step 1: annihilator}]
Take the dictionary $\{u_{tt},u_{xx},u_{tx},u_t,u_x,u\}$. The kernel of the Gram matrix is the line spanned by $u_{tt}-u_{xx}=0$. Output:
\[
P=\tau^2-\xi^2 .
\]
\end{example}

\begin{example}[{Step 2: factor}]
\[
P=Q_1Q_2,\qquad Q_1=\tau-\xi,\qquad Q_2=\tau+\xi .
\]
Two irreducible, simple factors.
\end{example}

\begin{example}[{Step 3: EP tests}]
\textbf{Remove $Q_1$.} We test $w_1=(\partial_t+\partial_x)u$. The two parts behave differently:
\[
(\partial_t+\partial_x)f(x+t)=2f'(x+t),\qquad (\partial_t+\partial_x)g(x-t)=0 .
\]
With $f(s)=1/(1+s^2)$:
\[
w_1=-\frac{4(x+t)}{\big(1+(x+t)^2\big)^2}.
\]
This is rational, not a finite sum of polynomials times exponentials: $w_1\notin\EP$.

\textbf{Remove $Q_2$.} Symmetrically, with $g(s)=e^{-s^2}$:
\[
w_2=(\partial_t-\partial_x)u=-2g'(x-t)=4(x-t)\,e^{-(x-t)^2}.
\]
A Gaussian times a polynomial is not an exponential-polynomial: $w_2\notin\EP$.
\end{example}

\begin{example}[{Step 4: reduce}]
Both tests say ``not EP'', so nothing is dropped. The point spectrum is empty, and the wave equation is identified:
\[
\Ann(u)=(\tau^2-\xi^2).
\]
\end{example}

\begin{example}[{Step 5: decompose}]
$(Q_1,Q_2)=(\tau,\xi)$, so $\delta(P)=1$. The components are $u_1=f(x+t)$ and $u_2=g(x-t)$, determined up to a constant moved from one to the other. Whether the components of every solution stay in $H^m$ on a given region is decided by Theorem~\ref{f:thm:hm}; on a square it holds (Example after Theorem~\ref{thm:density}).
\end{example}

\subsection{Worked example 2: advection--diffusion}

Now take the solution of Example~\ref{ex:advdiff},
\[
u=g(x-ct;\ \sigma^2+2Dt),\qquad c=2,\ D=\tfrac12,\ \sigma^2=1 .
\]

\begin{example}[{Steps 1--4}]
The Gram matrix of $\{u_t,u_x,u_{xx},u\}$ has kernel spanned by $u_t+2u_x-\tfrac12u_{xx}$, so
\[
P=\tau+2\xi-\tfrac12\xi^2 .
\]
It is irreducible (degree one in $\tau$). The test is on $u$ itself: a Gaussian is not an exponential-polynomial, so $\Ann(u)=(P)$ and the equation is identified.
\end{example}

\begin{example}[{Step 5}]
There is a single factor: no decomposition is needed.
\end{example}

\begin{example}[{Step 6: terms and interaction}]
The log-symbol is
\[
\log\hat u(t,\kappa)=\log\hat u_0(\kappa)+t\big(-2i\kappa-\tfrac12\kappa^2\big).
\]
Its two coefficients are the two terms: advection with speed $2$ and diffusion with constant $\tfrac12$. The interaction term follows from~\eqref{f:eq:eint}, at cost $O(2)$ per Fourier mode:
\[
\widehat{\Eint}(t,\kappa)=\big(e^{-2i\kappa t}-1\big)\big(e^{-\kappa^2t/2}-1\big)\,\hat u_0(\kappa).
\]
Figure~\ref{f:fig:eint} plots it.
\end{example}

\subsection{Cost and feasibility}

\begin{table}[!htb]
\centering\small
\begin{tabular}{@{}L{4.6cm}L{5.6cm}L{3.6cm}@{}}
\toprule
Component & Cost or feasibility & Source\\\midrule
annihilator from a dictionary & Gaussian elimination, $O(M^3)$ & Step 1\\
factorisation & standard computer algebra & Step 2\\
EP test, model M1 & syntactic & Step 3\\
EP test, model M3 & holonomic, decidable mod (AC) & Thm.~\ref{thm:holo}\\
EP test, black box & impossible without a bound & Thm.~\ref{f:thm:blackbox}\\
EP test with bound $B$ & linear algebra in degree $B+\deg P-1$ & Thm.~\ref{thm:decide}\\
components (evolution) & one rational multiplier per factor & Eq.~\eqref{f:eq:lagrange}\\
$\Eint$: value / by order / table & $O(N)$ / $O(N^2)$ / $\Theta(2^N)$ & Thm.~\ref{thm:cost}\\
\bottomrule
\end{tabular}
\caption{Cost and feasibility of the components of the algorithm.}\label{tab:complexity}
\end{table}

\begin{keypoint}
Every step of the algorithm is a theorem with explicit hypotheses. The algorithm is exact where the theorems are positive (Steps~2, 4, 5 and 6 on exponential-type data), and it reports honestly where they are negative (Step~3 on finite data, Step~6 in Sobolev spaces).
\end{keypoint}

